\documentclass[10pt]{article}

\usepackage[a4paper, total={6in, 9.5in}]{geometry}
\usepackage[dvipsnames,svgnames]{xcolor}
\usepackage[shortlabels]{enumitem}
\usepackage[framemethod=TikZ]{mdframed}
\usepackage{amsmath,amssymb,amsthm}
\usepackage[colorlinks]{hyperref}

\hypersetup{
	colorlinks=true,
	linkcolor=blue,
	filecolor=blue,      
	urlcolor=blue,
	pdftitle={DRP Notes}
}

\usepackage{mathtools}
\usepackage{thmtools}
\usepackage[textsize=scriptsize]{todonotes}

\renewcommand{\tilde}{\widetilde}
\renewcommand{\hat}{\widehat}
\renewcommand{\setminus}{\!-\!}
\providecommand{\ol}{\overline}
\providecommand{\eps}{\varepsilon}
\providecommand{\half}{\frac{1}{2}}
\providecommand{\dang}{\measuredangle} %% Directed angle
\providecommand{\dd}{\mathrm d}
\providecommand{\CC}{\mathbb C}
\providecommand{\FF}{\mathbb F}
\providecommand{\HH}{\mathbb H}
\providecommand{\NN}{\mathbb N}
\providecommand{\QQ}{\mathbb Q}
\providecommand{\RR}{\mathbb R}
\providecommand{\ZZ}{\mathbb Z}
\providecommand{\RP}{\mathbb RP}
\providecommand{\CP}{\mathbb CP}
\providecommand{\dg}{^\circ}
\providecommand{\ind}[1]{\mathbf 1_{#1}}
\providecommand{\ii}{\item}
\providecommand{\alert}[1]{{\sffamily\textbf{\textcolor{magenta}{#1}}}}
\providecommand{\opname}{\operatorname}
\providecommand{\li}{\opname{li}}
\providecommand{\ls}{\opname{ls}}
\providecommand{\supp}{\opname{supp}}
\providecommand{\im}{\opname{im}}
\providecommand{\Ext}{\opname{Ext}}
\providecommand{\Tor}{\opname{Tor}}
\providecommand{\inv}{^{-1}}
\providecommand{\setdiff}{\smallsetminus}
\providecommand{\mb}[1]{\mathbf{#1}}
\providecommand{\abs}[1]{\left\lvert{#1}\right\rvert}
\providecommand{\norm}[1]{\left\|{#1}\right\|}
\providecommand{\dif}{\;d}
\providecommand{\id}{\operatorname{id}}

\newcommand{\expow}{\Lambda}
\newcommand{\tensor}{\otimes}
\newcommand{\Hom}{\opname{Hom}}

\reversemarginpar
\setlength{\marginparwidth}{1in}

\mdfdefinestyle{mdbluebox}{roundcorner=10pt,innerbottommargin=9pt,
	linecolor=blue,backgroundcolor=TealBlue!5,}
\declaretheoremstyle[headfont=\sffamily\bfseries\color{MidnightBlue},
mdframed={style=mdbluebox},]{thmbluebox}

\mdfdefinestyle{mdredbox}{frametitlefont=\bfseries,innerbottommargin=8pt,
	nobreak=true,backgroundcolor=Salmon!5,linecolor=RawSienna,}
\declaretheoremstyle[headfont=\bfseries\color{RawSienna},
mdframed={style=mdredbox},headpunct={\\[3pt]},postheadspace=0pt,]{thmredbox}

\mdfdefinestyle{mdgreenbox}{linecolor=ForestGreen,backgroundcolor=ForestGreen!5,
	linewidth=2pt,rightline=false,leftline=true,topline=false,bottomline=false,}
\declaretheoremstyle[headfont=\bfseries\sffamily\color{ForestGreen!70!black},
mdframed={style=mdgreenbox},headpunct={ --- },]{thmgreenbox}

\mdfdefinestyle{mdorangebox}{linecolor=RedOrange,backgroundcolor=RedOrange!5,
	linewidth=2pt,rightline=false,leftline=true,topline=false,bottomline=false,}
\declaretheoremstyle[headfont=\bfseries\sffamily\color{RedOrange!70!black},
mdframed={style=mdorangebox},headpunct={ --- },]{thmorangebox}

\mdfdefinestyle{mdblackbox}{linecolor=black,backgroundcolor=RedViolet!5!gray!5,
	linewidth=3pt,rightline=false,leftline=true,topline=false,bottomline=false,}
\declaretheoremstyle[mdframed={style=mdblackbox}]{thmblackbox}

\declaretheoremstyle[spaceabove=6pt,spacebelow=6pt]{boldhead}

\declaretheorem[style=thmbluebox,name=Theorem,numberwithin=section]{theorem}

\declaretheorem[style=boldhead,name=Example,sibling=theorem]{example}

\declaretheorem[style=thmbluebox,name=Corollary,sibling=theorem]{corollary}
\declaretheorem[style=thmbluebox,name=Proposition,sibling=theorem]{prop}
\declaretheorem[style=thmbluebox,name=Lemma,sibling=theorem]{lemma}
\declaretheorem[style=thmbluebox,name=Theorem,numbered=no]{theorem*}
\declaretheorem[style=thmbluebox,name=Lemma,numbered=no]{lemma*}
\declaretheorem[style=thmgreenbox,name=Claim,numberwithin=theorem]{claim}
\declaretheorem[style=thmgreenbox,name=Claim,numbered=no]{claim*}
\declaretheorem[style=boldhead,name=Remark,sibling=theorem]{remark}
\declaretheorem[style=boldhead,name=Remark,numbered=no]{remark*}
\declaretheorem[style=thmgreenbox,name=Definition,sibling=theorem]{definition}
\declaretheorem[style=thmgreenbox,name=Definition,numbered=no]{definition*}
\declaretheorem[style=thmorangebox,name=Warning,sibling=theorem]{warning}
\declaretheorem[style=thmorangebox,name=Warning,numbered=no]{warning*}
\declaretheorem[style=thmorangebox,name=Note,numbered=no]{note}
\declaretheorem[style=thmblackbox,name=Exercise,sibling=theorem]{exercise}
\declaretheorem[style=thmblackbox,name=Exercise,numbered=no]{exercise*}

\newenvironment{soln}{\begin{proof}[Solution]}{\end{proof}}
\newenvironment{walkthrough}{\noindent \textbf{Walkthrough.}}{\medskip}
\newlist{walk}{enumerate}{3}
\setlist[walk]{label=\bfseries (\alph*)}

\providecommand{\pow}{\operatorname{Pow}}
%\providecommand{\vocab}[1]{{{\sffamily\textolor{ForestGreen}{#1}}}}
\providecommand{\vocab}[1]{{{\sffamily\textit{#1}}}}
\providecommand{\lcm}{\operatorname{lcm}}
\providecommand{\ord}{\operatorname{ord}}
\providecommand{\Int}{\operatorname{Int}}
\providecommand{\rk}{\operatorname{rk}}
\providecommand{\Mat}{\operatorname{Mat}}

\usepackage{asymptote}
\begin{asydef}
	defaultpen(fontsize(10pt));
	unitsize(1cm);
	size(8cm); // set a reasonable default
	usepackage("amsmath");
	usepackage("amssymb");
	settings.tex="pdflatex";
	settings.outformat="pdf";
	// Replacement for olympiad+cse5 which is not standard
	import geometry;
	// recalibrate fill and filldraw for conics
	void filldraw(picture pic = currentpicture, conic g, pen fillpen=defaultpen, pen drawpen=defaultpen)
	{ filldraw(pic, (path) g, fillpen, drawpen); }
	void fill(picture pic = currentpicture, conic g, pen p=defaultpen)
	{ filldraw(pic, (path) g, p); }
	// some geometry
	pair foot(pair P, pair A, pair B) { return foot(triangle(A,B,P).VC); }
	pair orthocenter(pair A, pair B, pair C) { return orthocentercenter(A,B,C); }
	pair centroid(pair A, pair B, pair C) { return (A+B+C)/3; }
	// cse5 abbreviations
	path CP(pair P, pair A) { return circle(P, abs(A-P)); }
	path CR(pair P, real r) { return circle(P, r); }
	pair IP(path p, path q) { return intersectionpoints(p,q)[0]; }
	pair OP(path p, path q) { return intersectionpoints(p,q)[1]; }
	path Line(pair A, pair B, real a=0.6, real b=a) { return (a*(A-B)+A)--(b*(B-A)+B); }
	// cse5 more useful functions
	picture CC() {
		picture p=rotate(0)*currentpicture;
		currentpicture.erase();
		return p;
	}
	pair MP(Label s, pair A, pair B = plain.S, pen p = defaultpen) {
		Label L = s;
		L.s = "$"+s.s+"$";
		label(L, A, B, p);
		return A;
	}
	pair Drawing(Label s = "", pair A, pair B = plain.S, pen p = defaultpen) {
		dot(MP(s, A, B, p), p);
		return A;
	}
	path Drawing(path g, pen p = defaultpen, arrowbar ar = None) {
		draw(g, p, ar);
		return g;
	}
\end{asydef}


\title{\vspace{-2.0cm}DRP Journal}
\author{\large Aditya Pahuja}
\date{\large \today}
\begin{document}

\maketitle
\tableofcontents
\pagebreak

\section{Rudiments of homotopy theory}
The $n$th homotopy group $\pi_n(X,x_0)$ is the set of homotopy classes
of maps $(I^n,\partial I^n)\to(X,x_0)$, where $I$ is the interval $[0,1]$.
Note that such maps correspond uniquely to maps $(S^n,s_0)\to(X,x_0)$
because $I^n/\partial I^n$ is homeomorphic to $S^n$.
This set of homotopy classes is made into a group as follows:
for maps $f,g\colon(I^n,\partial I^n)\to(X,x_0)$,
\begin{equation*}
    (f+g)(s_1,s_2,\dots,s_n) = \begin{cases}
	f(2s_1,s_2,\dots,s_n), & s_1\in[0,\half]\\
	g(2s_1-1,s_2,\dots,s_n), & s_1\in[\half,1].
    \end{cases}
\end{equation*}
The additive notation suggests that this group is abelian,
which is indeed true for $n > 1$:

\begin{prop}
    \label{prop:homotopy-groups-abelian}
    For $n\ge 2$, $\pi_n(X,x_0)$ is abelian.
\end{prop}

Additionally, $\pi_n$ is a functor: a map $\varphi\colon(X,x_0)\to(Y,y_0)$
induces the homomorphism
\[
    \varphi_*\colon\pi_n(X,x_0)\to\pi_n(Y,y_0),\qquad
    \varphi_*[f] = [\varphi\circ f].
\]
The higher homotopy groups satisfy some nice elementary properties
similar to those satisfied by the fundamental group:
\begin{prop}
    For all $n\ge 2$, the following properties hold.
    \label{prop:pi-n-props}
    \begin{enumerate}[(a)]
        \ii If $p\colon(\widetilde X,\widetilde x_0)\to(X,x_0)$
	is a covering map, $p_*\colon\pi_n(\widetilde X,\widetilde x_0)
	\to\pi_n(X,x_0)$ is an isomorphism.
	\ii If $(X_\alpha)_{\alpha\in A}$ is a set of topological spaces,
	there is an isomorphism
	$\pi_n(\prod_\alpha X_\alpha) \cong \prod_\alpha\pi_n(X_\alpha)$.
    \end{enumerate}
\end{prop}

If $X$ is path connected,
and $\gamma$ is a path from $x_1$ to $x_0$,
there is a change of basepoint isomorphism
$\beta_\gamma\colon\pi_n(X,x_0)\to\pi_n(X,x_1)$.
If we only consider loops $\gamma$ based at $x_0$,
then the map sending $[\gamma]\in\pi_1(X,x_0)$ to
$\beta_\gamma\in\operatorname{Aut}\pi_n(X,x_0)$ is
an action of $\pi_1$ on $\pi_n$.
For $n\ge 2$ this action makes $\pi_n$ into a
module over the group ring $\ZZ[\pi_1]$.

\begin{definition}[Relative homotopy groups]
    \label{def:relative-homotopy-group}
    Let $x_0\in A\subseteq X$ for spaces $A$, $X$.
    The \vocab{relative homotopy group} $\pi_n(X,A,x_0)$
    is the set of homotopy classes of maps
    $(D^n, S^{n-1}, s_0)\to(X,A,x_0)$
    under homotopies relative to $\partial S^{n-1}$, or, equivalently,
    maps $(I^n,\partial I^n,J^{n-1})\to(X,A,x_0)$
    under homotopies relative to $\partial I^n$, where $J^{n-1}$
    is the closure of the complement of a single face
    of $I^n$ in $\partial I^n$.
\end{definition}

Maps in the second sense can just as well be thought of as maps
$f\colon(I^n,\partial I^n)\to(X,A)$
such that $f(s_1,\dots,s_{n-1},0) = x_0$;
thus $\pi_n(X,A,x_0)$ can be made into a group for $n\ge 2$
(which is abelian for $n\ge 3$) by the usual composition operation.
Like with homology, these relative homotopy groups
fit into a long exact sequence.

\begin{prop}
    \label{prop:homotopy-long-exact-sequence}
    Let $(X,A,B)$ be a triple of spaces
    (so $B\subseteq A\subseteq X$) and $x_0\in B$.
    There is an exact sequence
    \begin{equation*}
	\cdots\to\pi_{n+1}(X,A,x_0)
	\stackrel\partial\to\pi_n(A,B,x_0)\stackrel{i_*}\to\pi_n(X,B,x_0)
	\stackrel{j_*}\to\pi_n(X,A,x_0)
	\stackrel\partial\to\pi_{n-1}(A,B,x_0)\to\cdots
    \end{equation*}
    where $i\colon(A,B,x_0)\hookrightarrow(X,B,x_0)$
    and $j\colon(X,B,x_0,x_0)\hookrightarrow(X,A,x_0)$
    are the inclusions and the connecting homomorphism
    $\partial$ is induced by the restriction
    of a map $(I^n,\partial I^n, J^{n-1})\to(X,A,x_0)$
    to a map $(I^{n-1},\partial I^{n-1})\to(A,x_0)$,
    with $I^{n-1}$ embedded in $I^n$ as the face
    that was deleted to make $J^{n-1}$.
\end{prop}

In the case where $B = x_0$ this gives an exact sequence
\begin{equation*}
    \cdots\to\pi_{n+1}(X,A,x_0)
    \stackrel\partial\to\pi_n(A,x_0)\stackrel{i_*}\to\pi_n(X,x_0)
    \stackrel{j_*}\to\pi_n(X,A,x_0)
    \stackrel\partial\to\pi_{n-1}(A,x_0)\to\cdots.
\end{equation*}

\begin{proof} 
    First we check exactness at $\pi_n(A,B,x_0)$.
    The composition $i_*\partial = 0$, since
    given continuous map $f\colon(I^{n+1},\partial I^{n+1},J^n)\to(X,A,x_0)$,
    the restriction $f|_{I^n}\colon(I^{n},\partial I^{n})\to(A,B)$
    is nullhomotopic relative to $\partial I^{n}$
    via $f$ itself because $f(J^{n}) = x_0$.
    To check that $\ker i_* = \im\partial$,
    suppose that $f\colon(I^{n},\partial I^{n},J^{n-1})\to(A,B,x_0)$
    is such that $i\circ f\colon(I^n,\partial I^n,J^{n-1})\to(X,B,x_0)$
    is nullhomotopic relative to $\partial I^n$.
    Then there is a homotopy of $i\circ f$ to the constant map equalling $x_0$,
    which yields a continuous map out of $I^{n+1}$
    such that $\partial I^{n+1}$ maps into $B$ and $J^n$ [].

    Next we check exactness at $\pi_n(X,B,x_0)$.
    \todo{finish}
\end{proof}

\todo{long exact sequence of a fibration}

\todo{fix definition of homotopy groups above}


\section{Differential forms}

\begin{definition}[Differential form]
    \label{def:differential-form}
    Let $M$ be an $n$-manifold and $\expow^k T^*M$
    the $k$th exterior power of the cotangent bundle $T^*M$.
    A \vocab{differential $k$-form}
    (often just called a \vocab{$k$-form}) is a section
    $\omega$ of $\expow^k T^*M$.
    The vector space of $k$-forms on $M$ is denoted
    $\Omega^k(M)$, and the graded algebra of all differential forms on $M$
    is denoted $\Omega^*(M)$.
\end{definition}

In local coordinates $(x_1,\dots,x_n)$, a $k$-form $\omega$
is expressible as
\begin{equation*}
    \sum_{1\le i_1<\cdots<i_k\le n} A_{i_1,\dots,i_k}
    \dd x_{i_1}\wedge\cdots\wedge\dd x_{i_k}.
\end{equation*}
where the $A_{i_1,\dots,i_k}\colon U\to\RR$
are smooth functions defined in the coordinate neighborhood.

The multiplication operation in $\Omega^*(M)$ is given by
just wedging the forms together; in local coordinates, this looks like
\begin{equation*}
    (f\dd x_{i_1}\wedge\cdots\wedge\dd x_{i_k})
    \wedge(g\dd x_{j_1}\wedge\cdots\wedge\dd x_{j_\ell})
    = fg\dd x_{i_1}\wedge\cdots\wedge\dd x_{i_k}
    \wedge\dd x_{j_1}\wedge\cdots\wedge\dd x_{j_\ell}
\end{equation*}
and extend by linearity for more complicated-looking forms.

The algebra of $k$-forms comes with an
\vocab{exterior differentiation} operator,
which makes it into a cochain complex
known as the \vocab{de Rham complex}.
\begin{prop}
    \label{prop:existence-uniqueness-exterior-derivative}
    There exists a unique endomorphism
    $\dd\colon\Omega^*(M)\to\Omega^*(M)$
    sending $k$-forms to $(k+1)$-forms such that
    \begin{enumerate}[(i)]
        \ii $\dd^2 = 0$;
	\ii $\dd(\omega\wedge\eta) = \dd\omega\wedge\eta
	+ (-1)^{\deg\omega}\omega\wedge\dd\eta$;
	\ii if $f\in C^\infty(M) = \Omega^0(M)$, then
	$\dd f\colon TM\to \RR$ is the differential of $f$.
    \end{enumerate}
\end{prop}

\begin{proof}
    Over open subsets of $\RR^n$, we can define
    \begin{equation*}
	\dd(f\dd x_{i_1}\wedge\cdots\wedge\dd x_{i_k})
	= \sum_{j=1}^n \frac{\partial f}{\partial x_j}\dd x_j\wedge
	\dd x_{i_1}\wedge\cdots\wedge\dd x_{i_k}
    \end{equation*}
    and extend by linearity;
    this is easily seen to satisfy the necessary properties.
    For an arbitrary $n$-manifolds, take a partition of unity
    subordinate to a covering of the manifold
    by coordinate neighborhoods; the properties in the proposition
    are well-behaved under linear combination, so
    summing over the partition of unity lets us define $\dd$ on any manifold.
    \todo{check independence of coordinate system}

    Now let $\dd$ and $\tilde\dd$ be two endomorphisms
    of $\Omega^*(M)$ with the given properties
    \todo{finish}
\end{proof}

\pagebreak
\section{Homology and cohomology}

A \vocab{simplex} of dimension $n$ is the simplest $n$-dimensional polytope,
analogous to how the triangle is the simplest polygon.
The canonical choice of simplex is
\begin{equation*}
    \Delta^n \coloneq \{(t_0,\dots,t_n) : 0\le t_i,\; t_0+\cdots+t_n\le 1\}.
\end{equation*}
the convex hull of the endpoints of the $n$ coordinate vectors.

\todo{write down definitions of homology cohomology etc
and many good geometric examples}

\subsection{Homological algebra and the universal coefficients theorems}

The goal is to algebraically relate integral homology,
homology with coefficients in a particular group $G$,
and cohomology with values in $G$.
This manifests as relating the homology of the singular chain complex
\begin{equation*}
    \dots\to C_{n-1}(X)\to C_n(X)\to C_{n+1}(X)\to\cdots
\end{equation*}
to that of the complexes
\begin{gather*}
    \dots\to C_{n+1}(X)\tensor G
    \to C_{n}(X)\tensor G
    \to C_{n-1}(X)\tensor G
    \to\dots
    \\
    \dots\to \Hom(C_{n-1}(X),G)
    \to \Hom(C_n(X),G)
    \to \Hom(C_{n+1}(X),G)
    \to\dots
\end{gather*}
(where $\otimes$ and $\Hom$ are over $\ZZ$).

\begin{prop}
    \label{prop:left-right-exactness}
    The functor $-\otimes G$ is right-exact;
    the functors $\Hom(-,G)$ and $\Hom(G,-)$ are left-exact.
\end{prop}

\begin{proof}
    Let $0\to A\overset u\to B\overset v\to C\to 0$
    be a short exact sequence of abelian groups.
    We can see that $v\tensor 1\colon B\tensor G\to C\tensor G$
    is surjective because given $c\in C$ and $b\in v\inv(c)$,
    $(v\tensor 1)(b\tensor g) = c\tensor g$.
    To get exactness at $B\tensor G$,
    we note that
    \begin{equation*}
	((v\tensor 1)\circ(u\tensor 1))(a\tensor g)
	= vu(a)\tensor g = 0
    \end{equation*}
    for all simple tensors $a\tensor g\in A\tensor G$,
    so $\ker(v\tensor 1)$ contains $\im(u\tensor 1)$.
    In the other direction, suppose that
    $\sum_i b_i\tensor g_i\in\ker(v\tensor 1)$.
    If $b_i = u(a_i)$, then $\sum_i b_i\tensor g_i
    = u\tensor 1(\sum_i a_i\tensor g_i)$,
    hence $\sum_i b_i\tensor g_i\in\im(u\tensor 1)$;
    thus $-\tensor G$ is right-exact.

    The proof for $\Hom$ is similar.
\end{proof}

\todo{ext and tor}

The main theorem is the following:
\begin{theorem}[Universal coefficients formulas]
    \label{thm:universal-coefficients}
    Let $X$ be a topological space, $n$ a nonnegative integer,
    and $G$ an abelian group. Then the sequences
    \begin{gather*}
	0\to H_n(X)\tensor G \to H_n(X ; G)\to\Tor_\ZZ(H_{n-1}(X),G)\to 0,\\
	0\to H^n(X;\ZZ)\tensor G \to H^n(X; G)
	\to\Tor_\ZZ(H^{n+1}(X;\ZZ),G)\to 0,\\
	0\to\Ext(H_{n-1}(X),G)\to H^n(X; G)\to\Hom(H_n(X),G)\to 0
    \end{gather*}
    are split exact.
\end{theorem}

\subsection{Products}











\end{document}
